\scanpage{271}
\mspage{138}

4. Cas des graphes triangulants.

Un graphe $(S,A_0)$ \inserted{sans boucles \deleted{isolés}} est dit \uline{triangulant} si

\[
\begin{aligned}
&\forall\,s,t\in S,\ \text{tels }\deleted{\unsure{\exists}}\text{que }\{s,t\}\in A_0,\quad
\exists\,u\in S\ \text{tel que }\{s,u\},\{t,u\}\in A_0\\
&\hspace{13em}\text{i.e tel que }\{s,t,u\}\in F_0
\end{aligned}\tag*{(1)}
\]

où on a posé

\[
F_0=\{f\in\unsure{\mathfrak P_3(S)|P_3(S)}\mid\forall\,a\in\mathfrak P_2(f),\ \text{on a }a\in A_0\}.\tag*{(2)}
\]

\deleted{La condition d'homogénéité \unsure{conve|comm}}
Lorsqu'on est en présence d'un \uline{multigraphe} à homothéties \deleted{de} $Z$, on dit qu'il est triangulant si chacune des structures de graphe est triangulante. Il suffit \inserted{pour ceci,} lorsque les \deleted{\unsure{ces}} structures sont isomorphes entre elles (``condition d'isomorphie'' du § 2) que l'une d'elles soit triangulante. \deleted{\unsure{4}} Mais il n'est \deleted{pas} sans doute pas vrai qu'il en soit ainsi dans tous les cas.
% archive sheet number: 270
